% ==================================================================
%  Infokommunikáció tankönyvek közös preambuluma (a tankönyvsablon alapján)
%  Fordítás: pdflatex, xelatex vagy lualatex (kétszer), illetve tectonic
% ==================================================================
\documentclass[a4paper,11pt,twoside,openany]{book}
\usepackage{iftex}
\ifPDFTeX
  \usepackage[T1]{fontenc}
  \usepackage[utf8]{inputenc}
  \usepackage{lmodern}
  \usepackage{textcomp}
\else
  \usepackage{fontspec}
\fi
\usepackage[magyar]{babel}
\usepackage{microtype}
\usepackage[a4paper,inner=24mm,outer=20mm,top=24mm,bottom=26mm]{geometry}
\usepackage{amsmath,amssymb,mathtools}
\usepackage{icomma}
\usepackage{xcolor}
\usepackage{tikz}
\usetikzlibrary{arrows.meta,calc,decorations.pathreplacing,positioning,patterns}
\usepackage{pgfplots}
\pgfplotsset{compat=1.18}
\usepackage{tabularx,booktabs,array,longtable}
\usepackage{enumitem}
\usepackage[most]{tcolorbox}
\usepackage{fancyhdr}
\usepackage{titlesec}
\usepackage{cancel}
\usepackage{adjustbox}
\usepackage[hidelinks]{hyperref}

% ---------- színek: csak a behelyettesítés jelzésére ----------
\definecolor{cA}{RGB}{29,95,196}    % kék
\definecolor{cB}{RGB}{192,84,0}     % narancs
\definecolor{cC}{RGB}{122,61,179}   % lila
\definecolor{cD}{RGB}{11,125,112}   % kékeszöld
\definecolor{cE}{RGB}{180,36,90}    % bíbor
\definecolor{cR}{RGB}{28,122,49}    % zöld: eredmény
\definecolor{keret}{gray}{0.35}
\definecolor{hatter}{gray}{0.95}
\newcommand{\ca}[1]{{\color{cA}#1}}
\newcommand{\cb}[1]{{\color{cB}#1}}
\newcommand{\cc}[1]{{\color{cC}#1}}
\newcommand{\cd}[1]{{\color{cD}#1}}
\newcommand{\ce}[1]{{\color{cE}#1}}
\newcommand{\res}[1]{{\color{cR}\boxed{#1}}}
\newcommand{\bh}[2]{\mbox{\(#1\;\longmapsto\;#2\)}\hspace{1.4em plus 0.5em}\allowbreak}
\newcommand{\behely}[1]{\par\smallskip{\raggedright\noindent\hangindent=2em\hangafter=1{\small\sffamily\textbf{Behelyettesítés:}}\hspace{0.8em}#1\par}\smallskip}
% a megoldási sor jellege (tétel, behelyettesítés, ...) és megjegyzése
\newcommand{\jegyzet}[2]{\par\vspace{-2pt}{\centering\footnotesize{\sffamily\scshape #1}\if\relax\detokenize{#2}\relax\else\quad #2\fi\par}\smallskip}

% ---------- matematikai rövidítések ----------
\newcommand{\ld}{\operatorname{ld}}
\newcommand{\arctg}{\operatorname{arctg}}
\newcommand{\tg}{\operatorname{tg}}
\newcommand{\sinc}{\operatorname{sinc}}
\newcommand{\F}{\mathcal{F}}
\newcommand{\eff}{_{\mathrm{eff}}}
\newcommand{\un}[1]{\,\mathrm{#1}}
\newcommand{\vect}[1]{\mathbf{#1}}
\renewcommand{\Re}{\operatorname{Re}}
\renewcommand{\Im}{\operatorname{Im}}

% ---------- dobozok ----------
\tcbset{base/.style={enhanced,breakable,sharp corners,boxrule=0.5pt,colframe=keret,colback=white,
  fonttitle=\sffamily\bfseries\small,coltitle=black,colbacktitle=hatter,left=3mm,right=3mm,top=1.5mm,bottom=1.5mm,
  before skip=8pt,after skip=8pt}}
\newtcolorbox{feladat}[1]{base,title={#1},colframe=black,boxrule=0.8pt}
\newtcolorbox{keresunk}{base,title={Mit keresünk?}}
\newtcolorbox{jeloles}{base,title={Jelölések ebben a feladatban}}
\newtcolorbox{elmelet}[1][Elméleti háttér]{base,title={#1}}
\newtcolorbox{eredmeny}{base,title={Eredmény},colframe=cR,colbacktitle=cR!10}
\newtcolorbox{figyelem}[1][Figyelem]{base,title={#1},colframe=cB,colbacktitle=cB!10}
\newtcolorbox{tipp}[1][Tipp]{base,title={#1},colframe=cA,colbacktitle=cA!8}
\newtcolorbox{rajz}[1][Rajz: ezt kell felrajzolni]{base,title={#1},colframe=cD,colbacktitle=cD!8}

% jelöléstáblázat egy feladathoz
\newenvironment{jeltabla}{\par\noindent\small\renewcommand{\arraystretch}{1.25}\tabularx{\dimexpr\linewidth-3pt\relax}{@{}>{\raggedright\arraybackslash}p{2.6cm}X>{\raggedright\arraybackslash}p{2.6cm}@{}}\toprule
\textbf{Jel} & \textbf{Jelentése ebben a feladatban} & \textbf{Érték, egység}\\\midrule}{\bottomrule\endtabularx\par}

% lépés
\newcounter{lepes}[section]
\newcommand{\lepes}[1]{\par\medskip\stepcounter{lepes}\noindent{\sffamily\bfseries\small\arabic{lepes}. lépés: #1}\par\nopagebreak\smallskip}

% ---------- TikZ: tollal rajzolható stílus ----------
\tikzset{
  toll/.style={line width=0.7pt,line cap=round,line join=round},
  vekony/.style={line width=0.4pt},
  szagg/.style={line width=0.5pt,dash pattern=on 3pt off 2pt},
  tengely/.style={line width=0.6pt,-{Stealth[length=2.2mm]}},
  nyil/.style={line width=0.6pt,-{Stealth[length=2mm]}},
  blokk/.style={draw,line width=0.7pt,minimum width=12mm,minimum height=8mm,font=\small},
}
\pgfplotsset{tollplot/.style={axis lines=middle,axis line style={line width=0.6pt,-{Stealth[length=2mm]}},
  every axis plot/.append style={line width=0.8pt,black},tick style={black},ticklabel style={font=\footnotesize},
  label style={font=\small},xlabel style={anchor=west},ylabel style={anchor=south},
  /pgf/number format/use comma,clip=false}}

% ---------- fejléc ----------
\pagestyle{fancy}
\fancyhf{}
\fancyhead[LE,RO]{\small\thepage}
\fancyhead[LO]{\small\sffamily\nouppercase{\leftmark}}
\renewcommand{\headrulewidth}{0.3pt}
\titleformat{\chapter}[display]{\normalfont\sffamily\bfseries}{\large\thechapter.\ \chaptertitlename}{4pt}{\LARGE}
\renewcommand{\chaptermark}[1]{\markboth{\thechapter.\ #1}{}}
\titlespacing*{\chapter}{0pt}{0pt}{18pt}
\titleformat{\section}{\normalfont\sffamily\bfseries\large}{\thesection}{0.8em}{}
\titleformat{\subsection}{\normalfont\sffamily\bfseries}{\thesubsection}{0.8em}{}
% a szakaszok a saját számukat (2.4, A1) a címükben hordozzák
\setcounter{secnumdepth}{0}
\setcounter{tocdepth}{1}
\setlist{itemsep=2pt,topsep=4pt}
\setlength{\parindent}{0pt}
\setlength{\parskip}{4pt}
\allowdisplaybreaks
\emergencystretch=2em
\raggedbottom
